\subsection{Consequence 8: Confinement and short-window persistence}
\begin{consequence}\label{cons:c8}
Under the declared shell-band extractor, shell signature, confinement predicate, and proximity criterion, a cited tetron-family window, including the locked tetrad alias when declared, may be reported as a confined short-window persistence class. Failure of the same predicate on the same shell-signature class triggers transmutation or recruitment thresholds on that window.
\end{consequence}
\paragraph{Depends on.} \Cref{choice:shell-band-extractor,def:shell-signature,def:confinement-predicate,def:closure-predicate,def:beat-cage-knot,def:cage-band,def:outer-enclosure-token,def:field-families,def:structural-constraints,cons:c4,cons:c6,cons:c7}.
\paragraph{Local substructure.} Shell-qualified tetron-family window; beat-knot / cage-knot typing; minimal cage-band $(4{:}6)$; outer enclosure token; $B^*$-dominance; shell-signature class; proximity criterion; confinement witness; transmutation threshold; torsional-load relation.
\paragraph{Derivation sketch.} On the cited motif-conditioned window, the closure predicate yields beat-knot and cage-knot classes on the canonical Q4 chart. The minimal cage-band and outer enclosure token supply the enclosure-side witness used by the shell-band extractor and shell signature. Shell qualification, dominant residual-duration class, and shell-signature stability then provide the persistence side of the report, while unresolved torsional load, recursive phase debt service, backlog forcing, and loss accumulation provide the failure side. The same declared proximity criterion separates sustained short-window persistence from transmutation / recruitment on that window family.
